% 例 3：将军饮马 — A,B 在直线 l 同侧，作 A' = A 关于 l 的对称点，连 A'B 交 l 于 P
\documentclass[tikz,border=4pt]{standalone}
\usepackage{ctex}
\usepackage{amsmath}
\usetikzlibrary{calc, decorations.markings}
\begin{document}
\begin{tikzpicture}[scale=0.9, thick,
  tickone/.style={postaction={decorate, decoration={markings,
    mark=at position 0.5 with {\draw (-1.5pt,-3pt) -- (1.5pt,3pt);}}}}]
  \draw (-3.5,0) -- (4.5,0) node[right]{$l$};
  \coordinate (A) at (-1.5, 1.5);
  \coordinate (B) at (2.5, 2.5);
  \coordinate (Ap) at (-1.5, -1.5);
  % 计算 P：在 A'B 与 l (y=0) 的交点
  % A'=(-1.5,-1.5)  B=(2.5,2.5) 直线方程 y - (-1.5)= ((2.5-(-1.5))/(2.5-(-1.5)))(x-(-1.5))= 1*(x+1.5)
  % y= x+0  令 y=0 → x=0
  \coordinate (P) at (0,0);

  % 垂线和对称
  \draw[gray, dashed] (A) -- (Ap);
  \draw (-1.7,0) -- (-1.7,0.2) -- (-1.5,0.2);
  \draw[tickone] (A) -- ($(A)!0.5!(Ap)$);
  \draw[tickone] ($(A)!0.5!(Ap)$) -- (Ap);

  % 路径 PA + PB（实线红）
  \draw[red, thick] (A) -- (P) -- (B);
  % A'B 蓝虚
  \draw[blue, dashed, thick] (Ap) -- (B);

  \fill (A) circle (1.4pt); \fill (B) circle (1.4pt);
  \fill (Ap) circle (1.4pt); \fill (P) circle (1.4pt);
  \node[above] at (A) {$A$};
  \node[above] at (B) {$B$};
  \node[below] at (Ap) {$A'$};
  \node[below right] at (P) {$P$};
\end{tikzpicture}
\end{document}
